\section{Ameaças à Validade}

Os pontos que mais podem impactar os resultados do presente são a escolha de
escopo, considerar funções apenas chamadas diretamente por outra no código, e os
projetos escolhidos.

O escopo foi escolhido de forma a mitigar possíveis erros enquanto controlando a
complexidade da tarefa de minerar as chamadas de funções do código. Apesar de
ser uma possível ameaça aos resultados, tal torna apenas os padrões encontrados
no presente estudo mais restritos do que talvez o seriam realmente. 

Um exemplo simples, porém comum, que escaparia a detecção dos métodos empregados
seria criar uma função que tracasse / liberasse um mutex e escrevesse algo no
canal de erro para fins de observabilidade. Em tal caso, as chamadas à funções
de trancar e liberar o mutex se encontrariam em funções distintas, mesmo que no
fluxo de execução do código estivessem no mesmo local.

Contudo, é importante apontar que, seguindo os presentes procedimentos, um
padrão que não ocorre jamais seria observado.

O dataset utilizado nas análises, os projetos escolhidos, também tem um grande
impacto nos padrões detectados. Supondo que os exemplos em si possuam erros,
como é sempre o caso para programas, os padrões de uso detectados nesses também
estariam errados. 

A fim de mitigar este problemas, os projetos para análise foram escolhidos
exatamente por sua popularidade e longevidade. Apesar de isto não garantir que
estejam livres de erros, espera-se que o fato de os programas em questão serem
utilizados por vários usuários indique minimamente que funcionem como esperado.
Dessa forma, como são projetos funcionais, espera-se que tais sigam minimamente
padrões de uso reais. 

%TODO: mencionar que padrões de uso != regras de uso
